\section{模运算与模逆}
\label{sec:04-Discrete-Mathematics/04-Number-Theory/02}

\subsection{一种你从小就见过但没认真对待过的运算}

现在是下午 2 点。80 小时后是几点？你不需要在日历上翻三天——80 除以 24 余 8，所以 80 小时后是晚上 10 点。这件事你小学就会做，但你可能没注意到你暗中使用了一个奇怪的代数规则：同一个``几点''可以对应无穷多个不同的整数——2 点、26 点、50 点指的都是针盘上同一个位置。

更怪的事在后面。普通算术里，两个非零的数乘起来一定是非零。但在钟面上：3 点钟方向出发，每次加 4 小时——从 3 到 7，再从 7 到 11。三次跳跃后，你回到了 3。用模运算写就是 \(3+4\times3\equiv3\pmod{12}\)，等价于 \(4\times3\equiv0\pmod{12}\)。两个非零的数（4 和 3）乘起来却在钟面上等于零。

这种``看似正常、实则不同''的运算规则，正是模运算一旦认真对待就会不断冒出来的惊喜。它看起来像普通代数，但``非零''不等于``可除''。一个余数类能否被约掉，唯一的门槛是它与模数是否互素。Bézout 恒等式既证明逆元存在，也实际构造它。

\subsection{跳格子：什么时候能走遍所有位置}

先看一个按固定格数跳动的钟面。模 \(n\) 的环上有 \(n\) 个位置。从某个刻度 \(i_0\) 出发，每次跨过 \(c\) 格，第 \(k\) 次到达

\[
i_k\equiv i_0+kc\pmod n.
\]

每次跳 4 格在模 12 下只会在三个刻度间循环：\(i_0, i_0+4, i_0+8\)，然后回到 \(i_0\)。每次跳 5 格却能在回到起点前经过全部十二个刻度。为什么？

因为你回到起点的那一刻，意味着 \(kc\) 是 \(n\) 的倍数——中间经过了 \(k\) 步，总共移动了 \(kc\) 格，落在了模 \(n\) 下的起点。第一个回到起点的步数是满足 \(n\mid kc\) 的最小 \(k>0\)，即

\[
k=\frac{n}{\gcd(c,n)}.
\]

若 \(\gcd(c,n)=1\)，则 \(k=n\)——你必须走遍全部 \(n\) 个位置才能回到起点。若 \(\gcd(c,n)=d>1\)，则只需 \(n/d\) 步就开始重复，只能访问 \(n/d\) 个不同的位置。

哈希表的线性探测沿用同一结构，只是钟面刻度换成了数组槽位。步长与表长互素时，探测序列遍历所有槽位；非互素步长会让部分槽位永远无法被检查——即使它们空着，你也会在已满的槽之间原地循环，然后宣告``表满了''。代数上，这正是``乘以 \(c\)''这个操作：只有当 \(c\) 存在模逆时，才会把余数类重新排列（permutation）而不是合并。

模逆何时存在、怎样求出，是一旦开始认真对待模运算就必须回答的两个问题。讨论要从同余的严格定义开始。

\subsection{同余把整数按余数分组}

固定 \(n\ge1\)，

\[
a\equiv b\pmod n
\]

表示 \(n\mid(a-b)\)。它是等价关系：自反（\(a\equiv a\)）、对称（\(a\equiv b\Rightarrow b\equiv a\)）、传递（\(a\equiv b,\ b\equiv c\Rightarrow a\equiv c\)）都可由整除定义直接验证——传递性只需把 \(a-b\) 与 \(b-c\) 相加。等价类就是余数类。模 \(n\) 运算只关心整数落在哪一类，不关心代表元具体多大。

同余与加法、减法、乘法相容：

\[
a\equiv b,\ c\equiv d\pmod n
\Longrightarrow
a+c\equiv b+d,\quad ac\equiv bd\pmod n.
\]

验证代入定义即可：写 \(b=a+kn,\ d=c+\ell n\)，展开后两边之差仍是 \(n\) 的倍数。这让你可以``算到一半取余''——计算 \(3^{100}\bmod 7\) 时不必先算出 \(3^{100}\) 这个 48 位数字，可以每一步平方后取余，把中间结果始终压在一个小范围内。这个简单的事实支撑了 RSA、Diffie--Hellman 和所有的模幂运算。

\subsection{约分并非总是合法：为什么不能随便消去}

普通代数里，从 \(ac=bc\) 且 \(c\neq0\) 总能推出 \(a=b\)。模运算里这条规则失效了。从

\[
ac\equiv bc\pmod n
\]

不能总推出 \(a\equiv b\pmod n\)。反例：\(2\cdot1\equiv2\cdot4\pmod6\)，但 \(1\not\equiv4\pmod6\)。

失败的原因写在等价形式里：\(ac\equiv bc\pmod n\) 等价于 \(n\mid c(a-b)\)。若 \(c\) 与 \(n\) 有公因子，\(n\) 的一部分``整除力''可以来自 \(c\) 而非来自 \((a-b)\)，因此 \(n\) 不必整除 \((a-b)\) 本身。只有当 \(\gcd(c,n)=1\) 时，\(n\) 的全部整除力必须落在 \((a-b)\) 上，消去律才安全，这正是上一节证明的结论：若 \(\gcd(c,n)=1\) 且 \(n\mid c(a-b)\)，则 \(n\mid(a-b)\)。

这也解释了刚才的``非零相乘得零''现象。\(3\times4\equiv0\pmod{12}\) 不是因为 3 或 4 在钟面上消失了，而是因为 3 和 4 各自携带了 12 的一部分素因子——3 贡献了 3，4 贡献了 4（即 \(2^2\)），合在一起凑齐了一个 12。没有与模数互素的``非零''，不是真正的``可逆元''，只是碰巧不等于余数零的那个类。

\subsection{模逆由 Bézout 恒等式判断并构造}

\(a\) 的模 \(n\) 逆是满足

\[
ax\equiv1\pmod n
\]

的 \(x\)。它存在，当且仅当 \(\gcd(a,n)=1\)。因为同余等价于存在整数 \(k\) 使得

\[
ax+nk=1,
\]

这正是 Bézout 方程。上一节的扩展 Euclid 算法不仅判断方程是否有解，还直接算出 \(x\) 和 \(k\)。

例如求 \(7^{-1}\pmod{26}\)。扩展 Euclid：

\[
\begin{aligned}
26 &= 3\cdot7 + 5\\
7  &= 1\cdot5 + 2\\
5  &= 2\cdot2 + 1\\
2  &= 2\cdot1 + 0
\end{aligned}
\]

回代：

\[
\begin{aligned}
1 &= 5 - 2\cdot2\\
  &= 5 - 2\cdot(7-1\cdot5) = 3\cdot5 - 2\cdot7\\
  &= 3\cdot(26-3\cdot7) - 2\cdot7 = 3\cdot26 - 11\cdot7.
\end{aligned}
\]

所以逆元为 \(-11\equiv15\pmod{26}\)，核对：\(7\cdot15=105=4\cdot26+1\equiv1\)。

Hill 密码的加密矩阵需要在模 26 下可逆——也就是行列式与 26 互素——原因就在这里。一个 2x2 矩阵 \(\begin{pmatrix}a&b\\c&d\end{pmatrix}\) 在模 26 下的逆需要 \((ad-bc)^{-1}\pmod{26}\) 存在，这要求行列式 \(\gcd(ad-bc,26)=1\)。

\subsection{解线性同余：gcd 告诉你答案个数}

方程

\[
ax\equiv b\pmod n
\]

有解，当且仅当 \(d=\gcd(a,n)\) 整除 \(b\)。必要性：若 \(x_0\) 是解，则存在 \(k\) 使 \(ax_0-b=kn\)，所以 \(d\mid(ax_0-kn)=b\)。充分性：若 \(d\mid b\)，方程两边同除以 \(d\)，化为在模 \(n/d\) 下求解 \(\frac{a}{d}x\equiv\frac{b}{d}\)，此时 \(\frac{a}{d}\) 与 \(\frac{n}{d}\) 互素，有唯一逆元，解模 \(n/d\) 唯一。

模 \(n\) 下的完整解集：若 \(d=1\)，解模 \(n\) 唯一。若 \(d>1\) 且可解，先求出模 \(n/d\) 下的唯一解 \(x_0\)，然后模 \(n\) 下所有解为

\[
x_0,\ x_0+n/d,\ x_0+2n/d,\ \ldots,\ x_0+(d-1)n/d,
\]

共 \(d\) 个不同的余数类。

不要在没有检查 gcd 之前就写 \(x\equiv a^{-1}b\)——符号 \(a^{-1}\) 的使用本身就预设了逆元存在。应先判断 \(d\mid b\) 是否成立，不可解时直接停在这里。

一个具体的例子：\(6x\equiv9\pmod{15}\)。\(\gcd(6,15)=3\)，而 \(3\mid9\)，可解。两边同除以 3：\(2x\equiv3\pmod5\)。\(2^{-1}\equiv3\pmod5\)，所以 \(x\equiv3\cdot3\equiv9\equiv4\pmod5\)。模 15 下的解集：\(4,\ 4+5=9,\ 4+10=14\)。可以逐个验证：\(6\cdot4=24\equiv9\)；\(6\cdot9=54\equiv9\)；\(6\cdot14=84\equiv9\)——三个都成立。

\subsection{环与域：合数模与素数模的分界线}

模 \(n\) 余数类连同加法与乘法构成环 \(\Z_n\)。当 \(n\) 是素数 \(p\) 时，每个非零类 1, 2, ..., \(p-1\) 都与 \(p\) 互素，因此都有乘法逆元——\(\Z_p\) 是有限域。合数模数会出现零因子，如 \(2\cdot3\equiv0\pmod6\)，其中 2 和 3 都不是零，但乘积却为零。

这一差别不是细枝末节。有限域上可以做线性代数，可以做多项式插值——因为域里除法的定义是完备的，方程组要么无解要么唯一解。环上则不然：你没法做 Gauss 消元的除法步骤，因为 pivot 可能没有逆元；即使有解，解也可能不唯一。这一差别直接决定了\hyperref[sec:04-Discrete-Mathematics/05-Coding-and-Polynomial-Methods/01]{``有限域与多项式插值''}及后续纠错码为何必须工作在有限域上：Reed--Solomon 码的编码和解码都需要在域中进行多项式求值和插值；若放在 \(\Z_{26}\) 这种环上，整个代数结构直接垮掉。

这也是密码学钟爱大素数或特定结构合数（如 RSA 的 \(n=pq\)）的核心原因——素数场里的可逆性让密钥生成和解密永远可行，而攻击者面对的却是模合数时的因子分解困难。
